% !TeX TS-program = xelatex

\documentclass{resume}
\ResumeName{田润泽}

\usepackage{graphicx}
\usepackage{tikz}

\begin{document}

\ResumeContacts{
  邮箱：\ResumeUrl{mailto:z5605012@ad.unsw.edu.au}{z5605012@ad.unsw.edu.au},%
  电话：+86 15632002968%
}

\begin{tikzpicture}[remember picture, overlay]
  \node [anchor=north east, inner sep=1cm] at (current page.north east)
    {\includegraphics[width=2.4cm]{image.png}};
\end{tikzpicture}

\ResumeTitle
\vspace{-0.8em}
\begin{center}
  籍贯：河北省邯郸市 \quad 出生年月：2001 年 10 月 \quad 政治面貌：共青团员
\end{center}
\vspace{-0.4em}

\section{教育背景}
\ResumeItem
[澳大利亚新南威尔士大学|计算机科学与技术硕士]
{澳大利亚新南威尔士大学（University of New South Wales）}
[\textnormal{计算机科学与技术硕士（人工智能方向）\quad GPA：84.6/100}]
[2024.09 -- 至今]

预计毕业时间：2026 年 8 月。主修课程：计算机网络与应用（89 分）、数据结构与算法（87 分）。

\ResumeItem
[云南大学|网络空间安全学士]
{云南大学（211/双一流大学）}
[\textnormal{网络空间安全学士\quad GPA：81.7/100}]
[2019.09 -- 2023.06]

主修课程：计算机网络（90 分）、操作系统与安全（90 分）。

荣誉奖项：2019--2020 学年二等奖学金；2020--2021 学年三等奖学金；2021--2022 学年二等奖学金；2023 年度优秀学生论文奖；第八届中国大学生“互联网+”创新创业大赛省级银奖。

\section{学术成果}
\begin{itemize}
  \item \textbf{田润泽}, 周宇龙, 朱洪, 等. 基于局部信息的服务迁移路径选择算法[J]. 计算机应用, 2024, 44(07):2168--2174.
  \item \textbf{关键词：}服务迁移路径；长短期记忆（LSTM）；神经网络；轨迹预测；道路网络匹配；动态路径规划。
  \item \textbf{论文状态：}二区第一作者，已发表。
  \item \textbf{计算机软件著作权：}基于深度学习的图像识别系统 V1.0 \hfill \textbf{登记号：}2023SR1150627
\end{itemize}

\section{科研经历}
\ResumeItem{\textbf{基于财报语义解析的企业可比性分析系统（CCF-A 区在投）}}
[项目负责人 | Euler AI | 金融科技项目]
[2025.06 -- 至今]

\textbf{导师：}新南威尔士大学计算机科学与工程学院 Zhengyi Yang 博士

\textbf{任务：}构建从公司财报中自动提取主营业务、进行语义解析和标准化的系统，用于企业间可比性分析。

\textbf{过程：}
\begin{enumerate}
  \item 利用 MinerU 实现财报文本的结构化抽取（业务结构、营收构成、经营分析、股权结构等）。
  \item 基于 LLM 的语义解析，识别主营业务关键词与语义占比，结合上下文与财务占比生成结构化数据。
  \item 构建标准化模块，通过关键词对齐与相似度计算（Word2Vec、Embedding、LLM 自动匹配）统一业务表达，并结合行业向量化 clustering 构建企业业务图谱。
\end{enumerate}

\textbf{成果：}成功实现从财报文本到结构化业务信息的自动化抽取，覆盖主营业务、营收构成和上下游关系。系统输出的业务标签对齐准确率达到 87\%，相较传统行业分类提升约 32\%。

\textbf{工具：}Python，PyTorch，MinerU，Word2Vec

\ResumeItem{\textbf{基于 Swin-Unet 的无人机图像枯死树分割}}
[项目负责人 | 新南威尔士大学 | 计算机视觉项目]
[2025.05 -- 2025.08]

\textbf{导师：}新南威尔士大学计算机科学与工程学院 Erik Meijering 博士

\textbf{任务：}设计并开发不同的计算机视觉方法来分割森林航拍图像中的枯死树。

\textbf{过程：}完成数据标注与切片，构建 RGB+NRG 多光谱融合管线，采用 Swin-UNet（移窗注意力 + U-Net 跳连）训练与验证，并以形态学/连通域过滤进行后处理，形成端到端分割流程。

\textbf{成果：}Foreground IoU 相比基准 UNet 提升 52.6\%，达到 0.3510；精确率 60.54\%，准确率 96.8\%，展现了 SwinUNet 在林业高精度监测中的潜力。

\textbf{工具：}PyTorch，OpenCV

\ResumeItem{\textbf{图像处理应用研究——虹膜识别系统改进}}
[项目负责人 | 新南威尔士大学 | 计算机视觉综合项目]
[2025.02 -- 2025.05]

\textbf{导师：}新南威尔士大学电气工程与电信学院 Siyuan Chen 博士

\textbf{任务：}通过机器学习方法提升虹膜识别系统在瞳孔扩张与收缩等复杂情况下的性能。

\textbf{过程：}构建完整虹膜识别系统并在传统方法上使用机器学习加以改进。分割部分包含传统方法（Canny 边缘检测 + Hough 霍夫变换）与 U-Net 分割对比；归一部分实现多尺度半径归一化；编码部分尝试两种基于掩码的编码（一维 Log-Gabor 滤波与局部二值模式 LBP）；匹配部分使用传统方法（汉明距离）与机器学习（ResNet-50 分类网络、Triplet 三元组网络）改进系统匹配精度。

\textbf{成果：}提交完整技术报告，提出的方案显著提升分割精度和匹配鲁棒性，对比同类算法匹配准确度领先 12.3\%，提高系统在真实环境（虹膜扩大、缩小）下的泛化能力。

\textbf{工具：}MATLAB

\section{项目经历}
\ResumeItem{\textbf{Walka：基于 Gemini 的智能文化旅行规划器}}
[\ResumeUrl{https://github.com/FFF108822/walkaTT}{github.com/FFF108822/walkaTT} | Hackathon 项目 | APAC Solution Challenge 2025]
[2025.02 -- 2025.06]

\textbf{成果：}实现了一个基于 Gemini 2.5 Pro 的旅行规划 agent，根据用户需求直接生成符合需求的旅行计划和行程地图。

\section{专业技能}
\begin{itemize}
  \item \textbf{外语水平：}英语：雅思 6.5 分，英语六级，熟练使用英语作为学术研究语言。
  \item \textbf{编程技能：}熟练掌握 Python, SQL, C, C++, MATLAB, HTML, CSS, JavaScript, TypeScript 等编程语言。
  \item \textbf{软件工具：}AI IDEs, n8n, PostgreSQL, MATLAB, Postman, Simulink, AutoCAD, Microsoft Office。
  \item \textbf{兴趣爱好：}乒乓球、徒步、游泳、滑雪。
\end{itemize}

\end{document}
